\section{Transformer-Based Models for Code Understanding and Generation}

%Itt is hiányzik a related work, meg pár hivatkozás, pl. a qwen coderekhez

The introduction of the Transformer architecture marked a major breakthrough in
sequence modeling and has become the foundation of modern approaches to code
understanding and generation. Unlike recurrent models, Transformers rely
entirely on self-attention mechanisms, enabling efficient parallel processing
and improved modeling of long-range dependencies \cite{vaswani2017}.

In recent years, several Transformer-based models have been specifically adapted
for programming languages. These models are typically pre-trained on large-scale
code corpora and are capable of capturing both syntactic and semantic aspects of
source code.

One of the earliest influential models is CodeBERT, a bimodal Transformer
trained on both natural language and programming language data. It learns joint
representations of code and text, making it suitable for tasks such as code
search and documentation generation \cite{feng2020codebert}.

GraphCodeBERT extends this approach by incorporating structural information from
code, particularly data flow. By modeling relationships between variables and
their dependencies, it captures deeper semantic information beyond token
sequences \cite{guo2020graphcodebert}.

Another important model is CodeT5, a unified encoder–decoder Transformer
designed for both code understanding and generation. It introduces
identifier-aware pre-training objectives, allowing the model to better capture
the semantics of variable names and improve performance across various tasks
\cite{wang2021codet5}.

More recent advances include large-scale generative models such as StarCoder and
Qwen Coder, which follow a decoder-only architecture similar to large language
models. These models are trained on massive code datasets and demonstrate strong
performance in code generation and completion tasks.

Despite their success, Transformer-based models still face limitations. In
particular, they often struggle with generalization across different datasets
and may fail to fully capture deeper semantic relationships in complex code
structures \cite{saberi2023model}.